\chapter{The Method of Characteristics}

\section{Derivation of the characteristic equations}

Please review how to use the method of characteristics to solve first-order linear PDEs up to first-order quasi-linear PDEs in MAT351.
We now start from the nonlinear equation which is the most general case.
Consider the boundary value problem
$$
\begin{cases}
F(Du,u,x)=0 & \text{in } U,\\
u=g & \text{on } \Gamma\subset \partial U,
\end{cases}
$$
where $U$ is open in $\mathbb{R}^n$.
The goal is to formulate an equivalent system of ODEs.

\subsection*{Setup}
Try to follow paths
$$
x(s)=\bigl(x^1(s),\ldots,x^n(s)\bigr), \qquad s\in I\subset \mathbb{R},
$$
so that we can follow the solution along the curve
$
s\mapsto x(s),
$
starting with initial point
$
x_0=x(0)\in \Gamma.
$
Assume $u$ is a $C^2$ solution to the PDE. Define
$$
z(s):=u(x(s)),
$$
and
$$
p(s):=Du(x(s))=\nabla u(x(s)).
$$
\begin{definition}
    The characteristic equations are defined as
    $$
    \begin{cases}
    \dot{p}(s)
    =
    -D_xF(p(s),z(s),x(s))
    -
    D_zF(p(s),z(s),x(s))p(s),
    \\[0.5em]
    \dot{z}(s)
    =
    D_pF(p(s),z(s),x(s))\cdot p(s),
    \\[0.5em]
    \dot{x}(s)
    =
    D_pF(p(s),z(s),x(s)).
    \end{cases}
    $$
\end{definition}

\begin{remark}
The three equations have the following meanings:
\begin{itemize}
    \item The equation for $\dot{x}$ chooses the curve $x(s)$ to move in the desired characteristic direction.
    \item The equation for $\dot{z}$ describes how $u(x(s))$ evolves with respect to the new time variable $s$.
    \item The equation for $\dot{p}$ comes from differentiating the PDE with respect to $x$. Only needed when the PDE is fully nonlinear.
\end{itemize}
\end{remark}

\subsection*{Derivation}
Start from the PDE
$
F(Du,u,x)=0.
$
Differentiate the PDE with respect to $x^i$. We denote the components of $F$ by
$
F=F(p,z,x).
$
Then
$$
\sum_{j=1}^n
\frac{\partial F}{\partial p^j}
\frac{\partial p^j}{\partial x^i}
+
\frac{\partial F}{\partial z}
\frac{\partial z}{\partial x^i}
+
\frac{\partial F}{\partial x^i}
=0.
$$
Equivalently,
\begin{equation}
\sum_{j=1}^n
F_{p^j}(Du,u,x)u_{x^i x^j}
+
F_z(Du,u,x)u_{x^i}
+
F_{x^i}(Du,u,x)
=0.
\tag{*}
\end{equation}
Now compute
$$
\dot{p}^i(s)
=
\sum_{j=1}^n
\frac{\partial p^i}{\partial x^j}
\frac{\partial x^j}{\partial s}.
$$
we get
\begin{equation}
\dot{p}^i(s)
=
\sum_{j=1}^n
u_{x^i x^j}(x(s))\dot{x}^j(s).
\tag{**}
\end{equation}
We want to match
$$
\sum_{j=1}^n
F_{p^j}(Du,u,x)u_{x^i x^j}
\text{ with }
\sum_{j=1}^n
u_{x^i x^j}(x(s))\dot{x}^j(s).
$$
So if we set
$
\dot{x}^j(s)=F_{p^j}(Du,u,x),
$
we have really nice structure. Intuitively, we choose the curve $x(s)$ to move in the direction
$
D_pF=F_p=(F_{p^1},\ldots,F_{p^n}).
$
Then equation $(**)$ becomes
$$
\dot{p}^i(s)
=
\sum_{j=1}^n
u_{x^i x^j}(x(s))F_{p^j}(Du,u,x).
$$
Plugging this into $(*)$, we get
$$
\dot{p}^i(s)
+
F_z(Du,u,x)u_{x^i}
+
F_{x^i}(Du,u,x)
=0.
$$



\section*{Example: First-order linear PDE}

\begin{claim}
For a first-order homogeneous linear PDE
$$
F(Du,u,x)=\langle b(x),Du(x)\rangle+c(x)u(x)=0,
$$
we can solve it by considering the characteristic equations
$$
\begin{cases}
\dot{x}(s)=b(x(s)),\\
\dot{z}(s)=-c(x(s))z(s).
\end{cases}
$$
\end{claim}

\begin{proof}
The function $F$ can be expressed as
$$
F(p,z,x)=b(x)\cdot p+c(x)z=0.
$$
Now consider the characteristic equations
$$
\begin{cases}
\dot{p}(s)
=
-D_xF(p(s),z(s),x(s))
-
D_zF(p(s),z(s),x(s))p(s),\\
\dot{z}(s)
=
D_pF(p(s),z(s),x(s))\cdot p(s),\\
\dot{x}(s)
=
D_pF(p(s),z(s),x(s)).
\end{cases}
$$
Starting from the equation for $\dot{x}$, since
$
D_pF=b,
$
we get
$
\dot{x}(s)=b(x(s)).
$
Next,
$
\dot{z}(s)=b(x(s))\cdot p(s).
$
Using the PDE
$
b(x)\cdot p+c(x)z=0,
$
we obtain
$
\dot{z}(s)=-c(x(s))z(s).
$
Notice that the equation for $\dot{p}$ is not needed here.
\end{proof}

\section{Mathematical technique: flattening the boundary}

\begin{definition}
Let $U\subset \mathbb{R}^n$ be open and bounded.
We say the boundary $\partial U$ is $C^k$ if, for every $x_0\in \partial U$, there exist $r>0$ and a $C^k$ change of coordinates
$$
\Phi:U\cap B(x_0,r)
\to
\left\{
x\in B(x_0,r): x^n>\sigma(x^1,\ldots,x^{n-1})
\right\}.
$$
Likewise, we say $\partial U$ is $C^\infty$ if this holds for $k=1,2,\ldots$.
We say $\partial U$ is analytic if the mapping $\sigma$ is analytic.
\end{definition}

\begin{definition}[Flattening the boundary]
The change of coordinates near a point of $\partial U$ is used to flatten out the boundary.
Fix $x_0\in \partial U$, and choose $r,\sigma$ as above. Define
$
y=\Phi(x)
$
by
$$
\begin{cases}
y^1=x^1,\\
\vdots\\
y^{n-1}=x^{n-1},\\
y^n=x^n-\sigma(x^1,\ldots,x^{n-1}).
\end{cases}
$$
Similarly, set the inverse
$
x=\Psi(y)
$
by
$$
\begin{cases}
x^1=y^1,\\
\vdots\\
x^{n-1}=y^{n-1},\\
x^n=y^n+\sigma(y^1,\ldots,y^{n-1}).
\end{cases}
$$
Then
$$
\Phi=\Psi^{-1},
\qquad
\det D\Phi=\det D\Psi=1.
$$
\end{definition}

\begin{claim}
After flattening the boundary, rewriting the original PDE for $u(x)$ as a new PDE for $v(y)$ has the form
$$
\begin{cases}
G(Dv,v,y)=0 & \text{in } V,\\
v=h & \text{on } \Delta.
\end{cases}
$$
\end{claim}

\begin{proof}
Suppose $u$ is a $C^1$ solution of the boundary value problem
$$
\begin{cases}
F(Du,u,x)=0 & \text{in } U,\\
u=g & \text{on } \Gamma\subset \partial U.
\end{cases}
$$
Rewrite
$
V:=\Phi(U),
$
and set
$
v(y):=u(\Psi(y)).
$
Equivalently,
$
u(x)=v(\Phi(x)).
$
By the chain rule,
$$
u_{x^i}(x)
=
\sum_{k=1}^n
v_{y^k}(\Phi(x))\Phi^k_{x^i}(x).
$$
In matrix form,
$$
D_xu(x)
=
D_yv(\Phi(x))D_x\Phi(x)
=
\nabla v(\Phi(x))^T D\Phi(x).
$$
Thus the PDE written in $y$ variables is
$$
F\bigl(Dv(y)D\Phi(\Psi(y)),v(y),\Psi(y)\bigr)=0.
$$
This is an expression of the form
$$
G(Dv(y),v(y),y)=0
\qquad \text{in } V.
$$
In addition,
$
v=h
\text{ on } \Delta,
$
where
$
\Delta=\Phi(\Gamma)
\text{ and }
h(y):=g(\Psi(y)).
$
\end{proof}

\section{Initial data for the characteristic ODE}

If we want to start a characteristic from the boundary point, the starting values cannot be chosen arbitrarily. They must match both the boundary condition and the PDE.
After straightening the boundary, we may assume that near
$
x_0\in \Gamma,
$
we have
$
\Gamma\subset \{x^n=0\}.
$
So locally the boundary is just the flat plane $x^n=0$.
Then the boundary condition
$$
u=g \quad \text{on } \Gamma
$$
becomes locally
$$
u(x^1,\ldots,x^{n-1},0)
=
g(x^1,\ldots,x^{n-1}).
$$
In the method of characteristics, one studies curves with unknowns
$$
x(s)\in \mathbb{R}^n,
\qquad
z(s)\in \mathbb{R},
\qquad
p(s)\in \mathbb{R}^n.
$$
At time $s=0$, the characteristic starts with
$$
x(0)=x_0,\qquad z(0)=z_0,\qquad p(0)=p_0.
$$
These are the initial conditions for the characteristic ODE, but not every triple
$
(p_0,z_0,x_0)
$
is acceptable.

\begin{definition}[Compatibility conditions]
The following conditions are called the \emph{compatibility conditions}:
$$
\begin{cases}
z_0=g(x_0),\\
p_0^i=g_{x^i}(x_0), \qquad i=1,\ldots,n-1,\\
F(p_0,z_0,x_0)=0.
\end{cases}
$$
A triple
$$
(p_0,z_0,x_0)\in \mathbb{R}^n\times \mathbb{R}\times \mathbb{R}^n
$$
verifying these conditions is called \emph{admissible}.
\end{definition}

\begin{remark}\
    \begin{itemize}
        \item The value $z_0$ is uniquely determined by the boundary condition and our choice of boundary point $x_0$.
        \item However, a vector $p_0$ satisfying the second and third compatibility conditions may not exist, or may not be unique.
    \end{itemize}
\end{remark}

\section{Noncharacteristic boundary condition}

If we have an admissible triple
$
(p_0,z_0,x_0),
$
can we perturb it?
For now, assume as above that
$
x_0\in \Gamma,
$
and that $\Gamma$ near $x_0$ lies in the flat plane
$
\{x^n=0\}.
$
Set
$
y=(y^1,\ldots,y^{n-1},0)\in \Gamma,
$
with $y$ close to $x_0$.

We want to solve the characteristic ODE with initial conditions
$$
p(0)=q(y),\qquad z(0)=g(y),\qquad x(0)=y.
$$
Our task now is to find a function $q(\cdot)$ such that
$
q(x_0)=p_0
$
and
$
(q(y),g(y),y)
$
is admissible.

\begin{lemma}[Noncharacteristic boundary conditions]
There exists a unique solution $q(\cdot)$ to the system
$$
\begin{cases}
q(x_0)=p_0,\\
q^i(y)=g_{x^i}(y), \qquad i=1,\ldots,n-1,\\
F(q(y),g(y),y)=0,
\end{cases}
$$
for all $y\in \Gamma$ close to $x_0$, provided that
$$
F_{p^n}(p_0,z_0,x_0)\neq 0.
$$
\end{lemma}

\begin{remark}\
    \begin{itemize}
        \item The condition
        $
        F_{p^n}(p_0,z_0,x_0)\neq 0
        $
        means that the characteristic curve starting from $x_0$ is not tangent to the flat boundary.
    \end{itemize}

\end{remark}

\begin{definition}[Noncharacteristic admissible triple]
In general, if $\Gamma$ is not flat near $x_0$, we say the admissible triple
$
(p_0,z_0,x_0)
$
is \emph{noncharacteristic} if
$$
D_pF(p_0,z_0,x_0)\cdot \nu(x_0)\neq 0,
$$
where $\nu(x_0)$ denotes the outward unit normal vector to $\Gamma$ at $x_0$.
\end{definition}
\begin{remark}\
    \begin{itemize}
        \item In the MAT351 notes, the PDE is more special:
        $$
        a(x,y,u)u_x+b(x,y,u)u_y=c(x,y,u).
        $$
        Equivalently,
        $$
        F(Du,u,x,y)
        =
        a(x,y,u)p_1+b(x,y,u)p_2-c(x,y,u)=0,
        $$
        where
        $
        p_1=u_x,
        $
        $
        p_2=u_y.
        $
        Then
        $
        F_{p_1}=a,$
        $
        F_{p_2}=b.
        $
        So the vector
        $$
        D_pF=(F_{p_1},F_{p_2})=(a,b)
        $$
        is the characteristic direction in the $(x,y)$-plane.
        
        For this quasilinear PDE, the characteristic ODE is
        $$
        \dot{x}=a(x,y,u),
        \qquad
        \dot{y}=b(x,y,u).
        $$
        So characteristics move in the direction
        $
        (a,b).
        $
        Now the boundary curve is parametrized by
        $
        \eta(t)=(k(t),l(t)).
        $
        Its tangent vector is
        $
        \eta'(t)=(\dot{k}(t),\dot{l}(t)).
        $

        The transversality condition says
        $$
        \det
        \begin{pmatrix}
        a(k(t),l(t),g(k(t),l(t))) & \dot{k}(t)\\
        b(k(t),l(t),g(k(t),l(t))) & \dot{l}(t)
        \end{pmatrix}
        \neq 0.
        $$
        This means the two vectors
        $
        (a,b)
        \text{ and }
        (\dot{k},\dot{l})
        $
        are not parallel.
        So geometrically, the characteristic direction is not tangent to the boundary.
        That is exactly what ``noncharacteristic boundary'' means.
    \end{itemize}

\end{remark}